\documentclass[11pt]{article}

\usepackage[margin=1in]{geometry}
\usepackage[T1]{fontenc}
\usepackage[utf8]{inputenc}
\usepackage{lmodern}
\usepackage{amsmath,amssymb,amsthm,mathtools}
\usepackage{booktabs,longtable,array}
\usepackage{tabularx}
\usepackage{enumitem}
\usepackage{hyperref}
\usepackage[numbers,square,sort&compress]{natbib}
\usepackage{microtype}

\hypersetup{
  colorlinks=true,
  linkcolor=blue,
  citecolor=blue,
  urlcolor=blue,
  pdftitle={A Residue-Validation and Adversarial-Control (RevA)-Gated Operator--Determinant Program for Prime-Induced Spectral Statistics},
  pdfauthor={Biju Jose}
}

\newcommand{\GEMLA}{$\Gamma$--EML--$\alpha$}
\newcommand{\EML}{\mathrm{EML}}
\newcommand{\ADL}{\mathrm{ADL}}
\newcommand{\RevA}{\textsc{RevA}}
\newcommand{\RevAFull}{Residue-Validation and Adversarial-Control}
\newcommand{\Rarith}{R_{\mathrm{arith}}}
\newcommand{\Kproj}{K_{\mathrm{proj}}}
\newcommand{\Keff}{K_{\mathrm{eff}}}
\newcommand{\XiF}{\Xi^{(F)}}
\newcommand{\Tr}{\operatorname{Tr}}
\newcommand{\rank}{\operatorname{rank}}
\newcommand{\diag}{\operatorname{diag}}
\newcommand{\argmin}{\operatorname*{arg\,min}}
\newcommand{\argmax}{\operatorname*{arg\,max}}
\newcommand{\dd}{\,\mathrm{d}}
\newcommand{\revtag}[1]{\textnormal{[Rev.~#1]}}

\newcolumntype{L}[1]{>{\raggedright\arraybackslash}p{#1}}
\newcolumntype{Y}{>{\raggedright\arraybackslash}X}

\theoremstyle{definition}
\newtheorem{definition}{Definition}
\newtheorem{remark}{Remark}
\newtheorem{protocol}{Protocol}
\newtheorem{gate}{Gate}
\theoremstyle{plain}
\newtheorem{proposition}{Proposition}
\newtheorem{lemma}{Lemma}
\newtheorem{theorem}{Theorem}

\title{A \RevAFull{} (\RevA)-Gated Operator--Determinant Program for Prime-Induced Spectral Statistics\\
\large From the Hybrid $\Gamma$--EML--$\alpha$ Architecture to Blind Zeta-Zero Non-Identification}
\author{
Biju Jose, Celes John, Syrita Maria Joe, Sincy Joe\\
Predilytics Inc., California, USA\\[0.4em]
{\footnotesize
\begin{tabular}{c}
\texttt{bjose@predilyticsinc.com}\\
\texttt{cjohnarrack3960@student.palomar.edu}\\
\texttt{joe010@csusm.edu}\\
\texttt{sj1188405@csuscd.edu}
\end{tabular}
}
}
\date{April 29 2026}

\begin{document}
\maketitle

\begin{abstract}
This paper develops a finite-dimensional, falsifiable operator--determinant program for prime-induced spectral statistics.  The construction is motivated by the Hilbert--Polya problem, by trace-formula analogies between primes and periodic orbits, and by the observed connection between zeta-zero statistics and random matrix theory.  It does not claim to prove the Riemann Hypothesis, to construct a Hilbert--Polya operator, or to identify the non-trivial zeros of the Riemann zeta function as a spectrum.  Instead, it isolates a sequence of experimentally testable operator conditions that a credible construction must satisfy before any such claim could be entertained.

The program combines three layers.  First, an Arithmetic Dynamical Layer (ADL) encodes prime powers as orbit lengths and gives an exact finite Euler-side trace/determinant identity.  Second, a Hybrid $\Gamma$--EML--$\alpha$ architecture supplies an action-based operator framework: the local action $\alpha$ defines a path measure, $\Gamma$ selects the induced minimum-action operator, and the EML representation supplies a common symbolic function class.  Third, a projection--residue determinant model separates a universal band-limited projection component from a trace-zero arithmetic residue.  A falsification gate, called \RevA{}, admits determinant phase-winding candidates only when the arithmetic residue is trace-zero, projection-orthogonal, aligned with the prime-dependent OLC residue, and separated from projection-only, random-residue, and RevC controls.

The numerical sequence reported here shows the following.  Phases \revtag{10.5--10.6} establish a stable projection--residue split.  Phases \revtag{10.7--10.8} produce the first persistent phase-winding tracks, but only under \RevA{} filtering.  The initial Phase \revtag{10.9} anti-overfit test fails under a corrected schema because the robust tracks are not \RevA-admissible.  Phases \revtag{10.9R--10.10R} repair this by evaluating RevA on the effective determinant residue $\gamma\cos(\phi)R$, after which out-of-sample and robustness tests pass with true-candidate RevA acceptance and false-control rejection.  Phases \revtag{10.11--10.13} then perform blind anchoring tests against known low zeta-zero ordinates.  The outcome is a controlled negative result: \RevA-admissible persistent determinant tracks survive structural repairs, but they do not anchor to zeta zeros and do not improve under $N$-refinement.  The central conclusion is therefore not an RH claim, but a methodological one: internal phase persistence is a real and reproducible operator phenomenon in the tested finite systems, while zeta-zero realization remains unachieved.
\end{abstract}

\paragraph{Keywords.} Riemann zeta function; Hilbert--Polya program; trace formula; sine kernel; Fredholm determinant; arithmetic dynamics; random matrix theory; operator learning; falsification.

\section{Introduction}

The Riemann Hypothesis is a statement about the non-trivial zeros of the Riemann zeta function, introduced in Riemann's 1859 memoir on the distribution of primes \citep{Riemann1859,Edwards1974,Titchmarsh1986}.  The famous spectral interpretation usually associated with the Hilbert--Polya idea asks whether those zeros can be realized as the spectrum of a self-adjoint operator.  Such a realization would explain the critical-line property through self-adjointness and would place the zeros in the same conceptual class as spectra of quantum Hamiltonians.  The difficulty is that no accepted operator with this property is known.  Random matrix theory supplies strong statistical evidence: Montgomery's pair-correlation work and Odlyzko's computations show that local zero statistics resemble unitary random-matrix statistics \citep{Montgomery1973,Odlyzko1987,Mehta2004}.  Keating and Snaith further connected zeta values on the critical line with characteristic polynomials of random unitary matrices \citep{KeatingSnaith2000}.  Yet statistical resemblance is not the same as a constructive spectral operator.

The present manuscript begins from the opposite direction.  Instead of assuming a spectral operator and searching for its relation to the zeros, we construct finite prime-induced operator systems, test them by trace and determinant identities, and only after passing falsification gates ask whether any zeta-zero behavior is visible.  The finite character of the construction is deliberate.  It makes every operator, determinant, trace, phase-winding count, and control experiment directly computable.  This is also the reason for the strict non-claim: the experiments do not establish the Riemann Hypothesis, do not identify zeta zeros, and do not provide a Hilbert--Polya operator.  They provide a disciplined path for determining which operator-level features are necessary, which are artifacts, and which survive out-of-sample tests.

The architectural motivation comes from the Hybrid $\Gamma$--EML--$\alpha$ framework \citep{JoseHybridGammaEMLAlpha2026}.  In that framework, $\alpha$ is a local action or path-weight contribution, $\Gamma$ is the deterministic minimum-action map induced by the probabilistic path measure, and EML is a symbolic representation class.  EML is used here as a compositional substrate rather than as a metaphysical claim: the recent EML operator construction shows that a single binary operation, together with constants, can generate a large class of elementary operations by uniform expression trees \citep{Odrzywolek2026}.  That representational discipline is useful in a program where phase features, logarithms, exponentials, and trigonometric terms must be generated without moving into an arbitrary black-box grammar.

Revision labels are shown throughout in bracketed form, for example 
\revtag{10.7}, 
\revtag{10.9R}, and 
\revtag{10.13z}.  These labels identify successive experimental revisions of the same operator--determinant program and should not be confused with numbered bibliographic citations, which are also rendered in square brackets by the citation style.

The paper is written as a research report rather than as a theorem asserting a final solution.  It consolidates the Phase 1--10.13 experimental program and reorganizes it into a submission-style argument: first lock the Euler side, then build a projection--residue determinant, then apply \RevA{} as an adversarial falsification gate, and only then perform zero-persistence and zeta-anchoring diagnostics.  The most important result is negative but constructive.  Persistent \RevA-admissible determinant phase tracks exist in the tested finite systems, but they remain distinct from zeta-zero realization.  This separation prevents a common overclaim: internal GUE-like or phase-winding behavior does not by itself imply the Riemann Hypothesis.

\section{Relation to earlier approaches}

Classical analytic approaches to the zeta function begin with analytic continuation, functional equations, and estimates for $\zeta(s)$ or related $L$-functions \citep{Edwards1974,Titchmarsh1986,Ivic1985}.  The operator program begins instead with Hilbert-space analogies, trace formulas, and spectral determinants.  Selberg's trace formula provides the paradigmatic setting in which spectral data and periodic orbit data are linked \citep{Selberg1956}, and quantum-chaotic trace formulas provide a parallel language of periodic orbits and spectral statistics \citep{Gutzwiller1990}.  Berry and Keating, Connes, and later authors proposed structural routes through quantization, scaling operators, or noncommutative geometry \citep{BerryKeating1999,Connes1999,Berry1985}.  These programs are mathematically rich, but they also illustrate the main obstruction: it is easier to describe what a Hilbert--Polya operator should do than to build one constructively from arithmetic primitives.

Random matrix theory occupies a different role.  It supplies a universal statistical model rather than an arithmetic operator.  The sine kernel appears as the local scaling limit of unitary invariant ensembles and governs the pair correlations of eigenvalues in the bulk \citep{Mehta2004,Forrester2010}.  The zeta zeros exhibit strikingly similar statistics after unfolding \citep{Montgomery1973,Odlyzko1987}.  For the present program, the lesson is not to insert the sine kernel by hand, but to force the structural properties whose unique translation-invariant band-limited limit is the sine kernel: projection, local unit density, translation invariance, compact spectral support, and near-binary Fourier multiplier.  If these conditions are learned from a $\Gamma$--EML--$\alpha$ action rather than imposed as a final kernel, the approach becomes less circular.

The program developed here differs from prior attempts in five ways.  First, it separates the Euler side from the zero side.  The Euler side must first be validated through finite prime-power traces and determinants; zero diagnostics are not allowed before that.  Second, it distinguishes global projection from local sine emergence.  A unitary conjugate $UPU^*$ is a projection automatically; this alone says nothing about sine-kernel emergence.  Third, it treats arithmetic as a residue coupled to a universal projection, rather than as something that should deform the projection itself.  Fourth, it uses adversarial controls through \RevA{}, so that projection-only and random-residue artifacts cannot be promoted as evidence.  Fifth, it treats non-identification as a valid scientific result.  If persistent determinant structures do not anchor to zeta zeros, the correct conclusion is failure of the current realization, not a weakened claim of success.

\section{Architecture: $\Gamma$--EML--$\alpha$ and the ADL layer}

The Hybrid $\Gamma$--EML--$\alpha$ architecture is an action-based framework over observed finite-dimensional state spaces.  In its general form, a local action $\alpha$ assigns a cost to a candidate transition, the cumulative action defines a path measure, and the deterministic operator $\Gamma$ selects the minimum-action transition.  If $X_t$ is the local context and $\phi'$ a candidate next state, an energy $\mathcal A(X_t,\phi')$ induces the Gibbs-type kernel
\begin{equation}
  p_{\mathcal A}(\phi'\mid X_t)
  = \frac{\exp[-\mathcal A(X_t,\phi')]}{Z_{\mathcal A}(X_t)},
  \qquad
  Z_{\mathcal A}(X_t)=\int \exp[-\mathcal A(X_t,z)]\,\dd \mu(z).
\end{equation}
The deterministic operator is then
\begin{equation}
  \Gamma_{\mathcal A}(X_t)=\argmin_{\phi'} \mathcal A(X_t,\phi'),
\end{equation}
and is a summary of the same probabilistic object rather than an unrelated map.  This matters because the experiments below are not merely matrix manipulations.  They are interpreted as finite operator objects selected by constraints, with falsification occurring at the level of observable traces, kernels, determinants, and phase behavior.

The Arithmetic Dynamical Layer supplies the arithmetic primitive.  Let
\begin{equation}
  \mathcal H_{P,K}=\operatorname{span}\{|p,k\rangle: p\le P,\;1\le k\le K\},
\end{equation}
where $p$ ranges over primes and $k$ over prime powers.  Define the diagonal length operator and weight operator
\begin{equation}
  A|p,k\rangle = k\log p\,|p,k\rangle,
  \qquad
  W|p,k\rangle = \log p\,|p,k\rangle.
\end{equation}
The weighted trace is
\begin{equation}
  L_{P,K}(s)=\Tr(We^{-sA})=
  \sum_{p\le P}\sum_{k\le K} \log(p) p^{-ks}.
\end{equation}
For $\Re(s)>1$, this finite trace approximates the logarithmic derivative $-\zeta'(s)/\zeta(s)$ as $P,K$ increase.  This construction is the Euler-side lock: before any zero-search object is considered, the system must reproduce prime-power trace behavior.

\begin{lemma}[finite Euler determinant]
Define
\begin{equation}
  \log D_{P,K}(s)=
  -\sum_{p\le P}\sum_{k\le K}\frac{1}{k}p^{-ks}.
\end{equation}
Then
\begin{equation}
  -\frac{\dd}{\dd s}\log D_{P,K}(s)=L_{P,K}(s).
\end{equation}
\end{lemma}

\begin{proof}
Differentiating term by term in the finite sum gives
\begin{equation}
  \frac{\dd}{\dd s}\left[-\frac{1}{k}p^{-ks}\right]
  = \log(p) p^{-ks}.
\end{equation}
Taking the negative derivative of the determinant logarithm yields the stated trace.  Since the sum is finite, no analytic-continuation issue enters this identity.
\end{proof}

This lemma explains why the determinant is weighted by $1/k$.  A naive determinant without that weight gives the wrong logarithmic derivative.  The trace--log identity is therefore not an aesthetic choice; it is the minimal finite determinant bridge to the Euler product.

\section{Projection emergence without inserting the sine kernel}

The sine kernel
\begin{equation}
  K_{\sin}(x,y)=\frac{\sin\pi(x-y)}{\pi(x-y)}
\end{equation}
is the local bulk correlation kernel associated with unitary random matrix ensembles after unfolding.  The weak version of a sine-kernel program is to insert this kernel after choosing unfolded coordinates.  The stronger version, used here, is to force the properties that make this kernel inevitable as a translation-invariant band-limited projection.  On $L^2(\mathbb R)$, a translation-invariant operator is diagonalized by Fourier transform.  If it is also a projection, its multiplier $m(\theta)$ satisfies $m^2=m$, so $m$ is an indicator function.  If the action enforces compact symmetric support $[-\Omega,\Omega]$, the inverse Fourier transform is
\begin{equation}
  K_\Omega(x,y)=\frac{1}{2\pi}\int_{-\Omega}^{\Omega}e^{i\theta(x-y)}\,\dd\theta
  =\frac{\sin\Omega(x-y)}{\pi(x-y)}.
\end{equation}
Thus the target is not ``put sine into the model.''  The target is ``learn or enforce projection, translation invariance, local density normalization, band limitation, and binary spectrum.''

The Phase 10 work used a sequence of projection repairs.  Soft fixed-point attempts reduced some defects but did not form a true projection algebra.  The later hard-projection refinements forced $K^2\approx K$ and then added Toeplitz and band-limit projections.  These steps confirmed that projection alone is insufficient.  A projection-only or unitary-conjugated projection can pass $K^2=K$ by construction while remaining far from the sine kernel.  The meaningful test is local/asymptotic: after restriction to central unfolded windows, does the kernel become Toeplitz, band-limited, density-stable, and close to the sine kernel without retuning?

This led to the central split
\begin{equation}
  \Keff = \Kproj + \lambda\Rarith,
\end{equation}
where $\Kproj$ is the fixed universal projection layer and $\Rarith$ is a trace-zero arithmetic residue.  The split is conceptually important.  The universal kernel must carry the sine/CLT structure.  The arithmetic information must enter as a controlled determinant-level residue, not by corrupting the universal projection.  This distinction was one of the key architectural lessons of Phases 10.5--10.6.

\section{The \RevAFull{} (\RevA) gate}

The \RevAFull{} (\RevA) gate was introduced to prevent false promotion of phase-winding artifacts.  In this paper, \RevA{} does not mean ``Revision A.''  It names a specific validation layer: residue validation of the effective arithmetic perturbation together with adversarial rejection of projection-only, random-residue, and RevC controls.  The need for such a gate became clear when projection-only and random-residue controls could produce raw determinant phase windings at high determinant scale.  A raw phase winding is therefore not meaningful unless its source residue passes arithmetic admissibility and the controls fail.

\begin{gate}[\RevAFull{} (\RevA) admissibility]
A determinant phase-winding candidate is \RevA-admissible only if the effective arithmetic residue $R_{\mathrm{eff}}$ satisfies all of the following conditions:
\begin{enumerate}[label=(\roman*)]
  \item trace neutrality: $|\Tr R_{\mathrm{eff}}|$ is below numerical tolerance;
  \item projection orthogonality: $|\langle R_{\mathrm{eff}},\Kproj\rangle|/(\|R_{\mathrm{eff}}\|_F\|\Kproj\|_F)$ is below tolerance;
  \item bounded norm: $\|R_{\mathrm{eff}}\|_F/\|\Kproj\|_F$ lies in a controlled interval;
  \item arithmetic alignment: $R_{\mathrm{eff}}$ has positive alignment with the prime-dependent OLC residue;
  \item adversarial rejection: projection-only, random trace-zero residue, and RevC controls do not pass the same gate.
\end{enumerate}
\end{gate}

The repaired form of \RevA{} evaluates
\begin{equation}
  R_{\mathrm{eff}}=\gamma\cos(\phi)R,
\end{equation}
not the raw residue $R$.  This correction matters because the successful determinant regime uses $\phi=\pi$, so the determinant actually sees a sign-flipped residue.  The original 10.9 audit rejected the true candidate by testing the wrong sign.  Phases \revtag{10.9R--10.10R} repaired this inconsistency and reran the out-of-sample tests under a single schema.

\section{Determinant observable and phase-winding tracks}

The determinant observable used in the finite tests has the form
\begin{equation}
  \Xi_{P}(s)=\det\left(I-\mu^2A_sA_{1-s}\right),
\end{equation}
where
\begin{equation}
  A_s=D_s\left(\Kproj+\epsilon\gamma\cos(\phi)\Rarith\right)D_{1-s}.
\end{equation}
Here $D_s$ is a diagonal scale operator depending on the chosen unfolded determinant coordinate, and $s=\sigma+it$ is sampled in a strip containing the critical line.  Candidate zeros are not accepted merely by small determinant magnitude.  They are identified by phase-winding behavior over small contours, tracked across prime cutoffs, and then filtered by persistence score and \RevA{}.

For a track $\tau$ observed at prime cutoffs $P_1<P_2<\cdots<P_m$, define the drift increments
\begin{equation}
  \Delta_j(\tau)=|s_\tau(P_{j+1})-s_\tau(P_j)|.
\end{equation}
The reported mean drift is the average of these increments.  A typical persistence score is
\begin{equation}
  S(\tau)=\frac{m}{M}\exp[-\overline{\Delta}(\tau)],
\end{equation}
where $M$ is the number of tested prime cutoffs.  A strong track requires high persistence, small drift, and repeated phase-winding hits.  The exact threshold is experimental, but the strict gate used in the later phases required $S>0.9$, at least one strong track, \RevA{} true pass, and \RevA{} control rejection.

\section{Experimental phase evolution}

The following table summarizes the phase progression in the operator--determinant program.  The values are finite-dimensional numerical diagnostics, not asymptotic theorems.  They are reported because they determine which claims are admissible and which are not.

\begin{longtable}{p{0.15\textwidth}p{0.37\textwidth}p{0.36\textwidth}}
\toprule
Phase & Main operation & Interpretation \\
\midrule
\revtag{10.5A--E} & Hard projection and arithmetic-retaining projection repairs & Projection algebra can be stabilized, but projection alone is not sine emergence. \\
\revtag{10.6} & Fixed $\Kproj$ plus trace-zero arithmetic residue & Structural pass; no zero-persistence claim. \\
\revtag{10.7} & Determinant scaling scan & First strong phase-winding persistence, but density drift remains high. \\
\revtag{10.8.1} & Local parameter refinement & Persistence preserved while raw density drift is reduced. \\
\revtag{10.8.2} & Density-observable stabilization & Density drift can be stabilized as an observable without creating the phase track. \\
\revtag{10.9.1--10.9.2} & Initial out-of-sample and robustness audit & Fails strict gate: robust-looking tracks are not \RevA-admissible. \\
\revtag{10.9R--10.10R} & Sign-aware RevA repair and multi-OOS robustness & Strict pass after evaluating the effective determinant residue. \\
\revtag{10.11} & Blind zeta-zero anchoring & \RevA-admissible persistent tracks do not anchor to zeta zeros. \\
\revtag{10.12a--10.12z} & Structural scale-law scan and blind anchoring & Persistence survives, but anchoring remains non-identification. \\
\revtag{10.13a--10.13z} & Riemann--von Mangoldt zeta-native scale laws & Smooth zero-density scaling does not improve anchoring. \\
\bottomrule
\end{longtable}

\section{Numerical results}

Table~\ref{tab:core-results} gives the core reported metrics from the late Phase 10 sequence.  The same cautious interpretation applies throughout: these are finite matrix tests on selected windows and cutoffs.  They are useful for falsification and prioritization, not for proof of RH.


\begin{table}[ht]
\centering
\footnotesize
\setlength{\tabcolsep}{3pt}
\renewcommand{\arraystretch}{1.12}
\begin{tabularx}{\linewidth}{@{}l L{0.18\linewidth} r r r Y@{}}
\toprule
Phase & Conditions & \shortstack{Strong\\tracks} & \shortstack{Best\\persistence} & \shortstack{Best\\mean drift} & Strict verdict \\
\midrule
\revtag{10.7} & scan & 1 & 0.9835 & 0.0167 & first persistence pass; high density drift \\
\revtag{10.8.1} & local refinement & 1 & 0.9835 & 0.0167 & pass: drift reduced \\
\revtag{10.9.1} & OOS audit & 0 & 0.2936 & 1.2257 & fail \\
\revtag{10.9.2} & robustness audit & 12 & 1.0000 & 0.0000 & fail: \RevA{} true pass false \\
\revtag{10.9R} & repaired OOS grid & 12 & 1.0000 & 0.0000 & pass \\
\revtag{10.10R} & multi-OOS robustness & 14 & 1.0000 & 0.0000 & pass \\
\revtag{10.11} & blind anchoring & 14 tested & -- & -- & non-identification \\
\revtag{10.12z} & blind anchoring & 8 tested & -- & -- & non-identification \\
\revtag{10.13z} & blind anchoring & 3 tested & -- & -- & non-identification \\
\bottomrule
\end{tabularx}
\caption{Core late-phase results.  The strict pass criterion after 10.9 requires strong persistence together with \RevA{} true-candidate acceptance and false-control rejection.  The anchoring phases test only strict-pass strong tracks against known zeta-zero ordinates.}
\label{tab:core-results}
\end{table}

The corrected interpretation of Phase 10.9 is particularly important.  In the initial \revtag{10.9} report, the robustness grid showed strong persistence-like tracks, but the true candidate failed \RevA{}.  Under the corrected schema, that is a fail, not a pass.  The repair in \revtag{10.9R} was not to lower the standard, but to correct the object being tested: RevA must evaluate the effective determinant residue.  After that repair, \revtag{10.9R} and \revtag{10.10R} pass the strict gate, with false controls still rejected.

The zeta-zero anchoring phases then deliberately attempt to falsify the interpretation that persistent tracks are zeta zeros.  In \revtag{10.11}, fourteen strict-pass strong tracks are tested.  There are no strong anchors within $\Delta t\le 0.5$, two weak proximities within $\Delta t\le 1.0$, and twelve nonmatches.  The best distance is about $0.5952$, and $N$-refinement does not improve.  In \revtag{10.12z}, structural scale laws are scanned blindly.  Persistence remains \RevA-admissible, but anchoring again fails: no promoted strong anchors, one weak proximity, seven nonmatches, and no $N$-refinement improvement.  In \revtag{10.13}, scale laws derived from the smooth Riemann--von Mangoldt term are tested.  CSV-level reconciliation shows that the best blind internal scale law is \texttt{rvm\_quantile}, not the baseline linear law.  The anchoring audit produces one raw strong proximity with best distance about $0.4775$, improving over the \revtag{10.12z} best raw distance, but the result is not promoted because $N$-refinement fails: the best and median distances worsen at larger $N$.

\section{Zeta-native scale laws and the 10.13 negative result}

The natural next idea after \revtag{10.12} is to ask whether the determinant coordinate should be scaled by the smooth zero-counting law rather than by an arbitrary internal coordinate.  The smooth Riemann--von Mangoldt term is
\begin{equation}
  \overline N(T)=\frac{T}{2\pi}\log\frac{T}{2\pi}-\frac{T}{2\pi}+\frac78.
\end{equation}
Phase \revtag{10.13} therefore tested structural scale laws derived from $\overline N(T)$: a count map, a quantile inverse, a cumulative density-weight map, and a hybrid.  None of these laws used individual zero ordinates during selection.  Selection remained blind and internal: persistence, \RevA{}, and density stability.

The result is informative precisely because it is mixed but ultimately negative.  The smooth zero-density law improves one raw proximity, but it does not solve the anchoring problem.  This suggests that the missing ingredient is not merely the average density of zeros.  If a determinant model of this type is to approach the zeta zeros, it likely needs the oscillatory explicit-formula correction that couples primes and zeros, not only the smooth Riemann--von Mangoldt density.  In other words, the program has separated three phenomena: universal projection persistence, arithmetic residue persistence, and actual zeta-zero anchoring.  Only the first two are currently observed.

\section{What differentiates this program}

The first differentiator is non-circularity.  The sine kernel is treated as a diagnostic endpoint of projection, translation invariance, local density, and band-limit constraints.  It is not used as the original definition of the arithmetic kernel.  The model can therefore fail to approach the sine kernel, and several intermediate versions did fail.  That failure information is central to the design.

The second differentiator is the separation of universal and arithmetic layers.  In many informal Hilbert--Polya attempts, arithmetic is pushed directly into a candidate Hamiltonian or into a spectral list.  Here the arithmetic layer is first locked by finite Euler trace/determinant identities.  The universal projection layer is then stabilized separately.  Arithmetic enters the determinant through a trace-zero residue, and that residue is subjected to \RevA{}.

The third differentiator is adversarial validation.  \RevA{} is not a cosmetic diagnostic; it changes the logical status of the results.  Raw phase windings are known to occur in controls.  The only admissible evidence is phase-winding persistence that survives out-of-sample tracking while the true effective residue passes and the controls fail.  This gives the program a built-in mechanism for rejecting fake GUE or fake zero emergence.

The fourth differentiator is the willingness to report non-identification.  The most scientifically important result of Phases 10.11--10.13 is not that the model almost finds zeros, but that it does not.  The determinant phase tracks are persistent and \RevA-admissible, but they do not refine toward known zeta-zero ordinates.  That conclusion prevents the framework from becoming another numerology exercise.  It also identifies the next obstruction: the oscillatory explicit-formula correction is not yet represented in the determinant scale law.

\section{Limitations and non-claims}

The main limitation is finite dimensionality.  All reported experiments use finite prime cutoffs, finite matrix windows, and finite determinant grids.  Their purpose is to test mechanisms and falsify artifacts, not to establish asymptotic theorems.  A second limitation is the dependence on unfolding and determinant coordinate choices.  Phases 10.12--10.13 show that structural scale laws can preserve internal persistence without producing zeta anchoring.  This means that scale-law selection remains an unresolved issue.

The strongest non-claim is that no result in this paper proves the Riemann Hypothesis.  No result identifies the non-trivial zeros of $\zeta(s)$ as the spectrum of a constructed self-adjoint operator.  No result establishes a Hilbert--Polya realization.  The phrase ``zeta-native'' in Phase 10.13 refers only to scale laws derived from the smooth Riemann--von Mangoldt term; it does not mean that the resulting determinant is the completed zeta function.

A further limitation concerns \RevA{} itself.  \RevA{} is a powerful falsification gate for the experiments reported here, but it is not a theorem of arithmetic.  It encodes structural requirements that were empirically necessary to reject known controls.  Future work should formalize \RevA{} as a proposition about invariant subspaces, projection-orthogonal residues, and determinant phase stability.  Until then, \RevA{} should be described as an experimental gate, not a mathematical proof of authenticity.

\section{Reproducibility and experimental protocol}

A reproducible version of the program should retain the following order.  First, build the finite prime-power ADL operator and validate the weighted Euler trace.  Second, validate the finite determinant and log-derivative identity.  Third, construct or learn a projection kernel and test projection, Toeplitz defect, band limitation, eigenvalue collapse, and local sine distance.  Fourth, split the effective kernel into $\Kproj$ and $\Rarith$, and verify trace-zero, projection-orthogonality, norm control, and OLC alignment.  Fifth, run determinant phase-winding searches.  Sixth, apply \RevA{} controls.  Seventh, track persistent candidates across prime cutoffs and robustness grids.  Eighth, and only then, perform zeta-zero anchoring.

The anchoring protocol must remain blind.  Known zero ordinates should not be used to select scale laws, tune determinant parameters, choose windows, or filter candidates.  They should enter only after strict-pass tracks have been selected by internal criteria.  A promoted anchor should require proximity to a known zero, persistence across cutoffs, \RevA{} admissibility, rejection of controls, and improvement under refinement.  The current experiments fail this final promotion criterion.

\section{Outlook}

The next logical step is an oscillatory explicit-formula correction that is still blind to zero ordinates.  The smooth Riemann--von Mangoldt term captures average zero density, but it omits the oscillatory prime-zero coupling that appears in explicit formulae.  A non-circular version of the next phase would use only prime-side oscillatory data to perturb the determinant coordinate or residue, then rerun the same \RevA{} and blind anchoring gates.  Success would require not only better proximity, but monotone improvement under $P$ and $N$ refinement without retuning.

A second direction is theorem-building around the projection--residue split.  One can try to prove that a trace-zero, projection-orthogonal residue perturbs determinant phase without destroying the band-limited projection structure, at least under finite-rank assumptions.  A third direction is to formalize \RevA{} as an invariant test.  If \RevA{} can be stated as a finite-dimensional theorem about admissible arithmetic residues and control rejection, it would elevate the program from a numerical pipeline to a theorem-guided experimental architecture.

\section{Conclusion}

This paper has organized the $\Gamma$--EML--$\alpha$ and ADL experimental program into a disciplined operator--determinant framework.  The positive result is that finite prime-induced systems can produce \RevA-admissible persistent determinant phase tracks after projection--residue separation and effective-residue repair.  The negative result is that these tracks do not anchor to zeta zeros under blind tests, structural scale-law repairs, or smooth Riemann--von Mangoldt scaling.

The conclusion is therefore deliberately restrained.  The program has not realized Hilbert--Polya and has not proved RH.  It has, however, produced a falsifiable pathway that separates universal projection structure, arithmetic residue structure, determinant phase persistence, and zeta-zero anchoring.  That separation is valuable.  It makes clear that internal phase persistence is not enough, and it identifies the missing mathematical ingredient: a non-circular representation of the oscillatory explicit-formula coupling between primes and zeros.

\appendix

\section{Finite-dimensional definitions used in the experiments}

\begin{definition}[prime-power orbit length]
For a prime $p$ and exponent $k\ge1$, the orbit length is
\begin{equation}
  x_{p,k}=k\log p.
\end{equation}
The finite ADL Hilbert space $\mathcal H_{P,K}$ is spanned by the corresponding basis vectors $|p,k\rangle$.
\end{definition}

\begin{definition}[phase-winding candidate]
A candidate is a grid cell in the $s=\sigma+it$ strip for which the determinant phase changes by approximately an integer multiple of $2\pi$ around a small contour, together with a local determinant magnitude diagnostic.  A candidate becomes a track only if it can be matched across multiple prime cutoffs.
\end{definition}

\begin{definition}[strong track]
A track is called strong in the reported experiments if it has high persistence score, low mean drift, and multiple phase-winding hits.  In the late strict gate, a phase is promoted only if at least one strong track exists and \RevA{} accepts the true effective residue while rejecting controls.
\end{definition}

\section{Bibliographic note}

The references below include classical analytic-number-theory sources, random-matrix and trace-formula sources, determinant/operator references, and internal architecture documents.  The internal documents are included to document the construction and experimental sequence; they should be replaced by archived code, data, and public preprints in any formal external submission.

\begin{thebibliography}{23}
\providecommand{\natexlab}[1]{#1}
\providecommand{\url}[1]{\texttt{#1}}
\expandafter\ifx\csname urlstyle\endcsname\relax
  \providecommand{\doi}[1]{doi: #1}\else
  \providecommand{\doi}{doi: \begingroup \urlstyle{rm}\Url}\fi

\bibitem[Riemann(1859)]{Riemann1859}
Bernhard Riemann.
\newblock {Ueber die Anzahl der Primzahlen unter einer gegebenen Gr\"osse}.
\newblock \emph{Monatsberichte der K\"oniglich Preu\ss ischen Akademie der
  Wissenschaften zu Berlin}, pages 671--680, 1859.
\newblock English translation by David R. Wilkins available online.

\bibitem[Edwards(1974)]{Edwards1974}
Harold~M. Edwards.
\newblock \emph{Riemann's Zeta Function}.
\newblock Academic Press, New York, 1974.

\bibitem[Titchmarsh(1986)]{Titchmarsh1986}
E.~C. Titchmarsh.
\newblock \emph{The Theory of the Riemann Zeta-Function}.
\newblock Oxford University Press, 2 edition, 1986.
\newblock Revised by D. R. Heath-Brown.

\bibitem[Montgomery(1973)]{Montgomery1973}
Hugh~L. Montgomery.
\newblock The pair correlation of zeros of the zeta function.
\newblock In \emph{Analytic Number Theory}, volume~24 of \emph{Proceedings of
  Symposia in Pure Mathematics}, pages 181--193, Providence, RI, 1973. American
  Mathematical Society.

\bibitem[Odlyzko(1987)]{Odlyzko1987}
Andrew~M. Odlyzko.
\newblock On the distribution of spacings between zeros of the zeta function.
\newblock \emph{Mathematics of Computation}, 48\penalty0 (177):\penalty0
  273--308, 1987.
\newblock \doi{10.1090/S0025-5718-1987-0866115-0}.

\bibitem[Mehta(2004)]{Mehta2004}
Madan~Lal Mehta.
\newblock \emph{Random Matrices}.
\newblock Elsevier Academic Press, 3 edition, 2004.

\bibitem[Keating and Snaith(2000)]{KeatingSnaith2000}
Jon~P. Keating and Nina~C. Snaith.
\newblock Random matrix theory and $\zeta(1/2+it)$.
\newblock \emph{Communications in Mathematical Physics}, 214:\penalty0 57--89,
  2000.
\newblock \doi{10.1007/s002200000261}.

\bibitem[Jose(2026{\natexlab{a}})]{JoseHybridGammaEMLAlpha2026}
Biju Jose.
\newblock Hybrid $\gamma$--eml--$\alpha$ architecture: A framework on
  multi-scale dynamics, action-based inference, and observable path measures,
  April 2026{\natexlab{a}}.
\newblock Internal manuscript draft, Predilytics Inc.

\bibitem[Odrzywolek(2026)]{Odrzywolek2026}
Andrzej Odrzywolek.
\newblock All elementary functions from a single binary operator, 2026.
\newblock Version 2, April 2026.

\bibitem[Ivi\'{c}(1985)]{Ivic1985}
Aleksandar Ivi\'{c}.
\newblock \emph{The Riemann Zeta-Function: Theory and Applications}.
\newblock John Wiley \& Sons, New York, 1985.

\bibitem[Selberg(1956)]{Selberg1956}
Atle Selberg.
\newblock Harmonic analysis and discontinuous groups in weakly symmetric
  riemannian spaces with applications to dirichlet series.
\newblock \emph{Journal of the Indian Mathematical Society}, 20:\penalty0
  47--87, 1956.

\bibitem[Gutzwiller(1990)]{Gutzwiller1990}
Martin~C. Gutzwiller.
\newblock \emph{Chaos in Classical and Quantum Mechanics}.
\newblock Springer, 1990.

\bibitem[Berry and Keating(1999)]{BerryKeating1999}
Michael~V. Berry and Jon~P. Keating.
\newblock $h=xp$ and the riemann zeros.
\newblock \emph{Supersymmetry and Trace Formulae: Chaos and Disorder}, pages
  355--367, 1999.

\bibitem[Connes(1999)]{Connes1999}
Alain Connes.
\newblock Trace formula in noncommutative geometry and the zeros of the riemann
  zeta function.
\newblock \emph{Selecta Mathematica}, 5:\penalty0 29--106, 1999.
\newblock \doi{10.1007/s000290050042}.

\bibitem[Berry(1985)]{Berry1985}
Michael~V. Berry.
\newblock Semiclassical theory of spectral rigidity.
\newblock \emph{Proceedings of the Royal Society of London. Series A},
  400\penalty0 (1819):\penalty0 229--251, 1985.
\newblock \doi{10.1098/rspa.1985.0078}.

\bibitem[Forrester(2010)]{Forrester2010}
Peter~J. Forrester.
\newblock \emph{Log-Gases and Random Matrices}.
\newblock Princeton University Press, 2010.

\bibitem[Katz and Sarnak(1999)]{KatzSarnak1999}
Nicholas~M. Katz and Peter Sarnak.
\newblock \emph{Random Matrices, Frobenius Eigenvalues, and Monodromy},
  volume~45 of \emph{Colloquium Publications}.
\newblock American Mathematical Society, 1999.

\bibitem[Bohigas et~al.(1984)Bohigas, Giannoni, and
  Schmit]{BohigasGiannoniSchmit1984}
Oriol Bohigas, Marie-Joya Giannoni, and Charles Schmit.
\newblock Characterization of chaotic quantum spectra and universality of level
  fluctuation laws.
\newblock \emph{Physical Review Letters}, 52\penalty0 (1):\penalty0 1--4, 1984.
\newblock \doi{10.1103/PhysRevLett.52.1}.

\bibitem[Costin and Lebowitz(1995)]{CostinLebowitz1995}
Ovidiu Costin and Joel~L. Lebowitz.
\newblock Gaussian fluctuation in random matrices.
\newblock \emph{Physical Review Letters}, 75\penalty0 (1):\penalty0 69--72,
  1995.
\newblock \doi{10.1103/PhysRevLett.75.69}.

\bibitem[Soshnikov(2000)]{Soshnikov2000}
Alexander Soshnikov.
\newblock Determinantal random point fields.
\newblock \emph{Russian Mathematical Surveys}, 55\penalty0 (5):\penalty0
  923--975, 2000.
\newblock \doi{10.1070/RM2000v055n05ABEH000321}.

\bibitem[Simon(2005)]{Simon2005}
Barry Simon.
\newblock \emph{Trace Ideals and Their Applications}.
\newblock American Mathematical Society, 2 edition, 2005.

\bibitem[Deift(1999)]{Deift1999}
Percy Deift.
\newblock \emph{Orthogonal Polynomials and Random Matrices: A Riemann--Hilbert
  Approach}.
\newblock American Mathematical Society, 1999.

\bibitem[Jose(2026{\natexlab{b}})]{JosePhase10Corpus2026}
Biju Jose.
\newblock Phase 10 experimental corpus for a reva-gated prime-induced
  operator--determinant program, April 2026{\natexlab{b}}.
\newblock Internal experimental reports and result tables for Phases
  10.3--10.13.

\end{thebibliography}


\end{document}